\section{Open Problems and Research Roadmap}
\label{sec:open}

The survey presented in the preceding sections makes clear that the field of
\nvs{}-to-world-modeling is progressing rapidly along multiple fronts simultaneously.
Per-scene reconstruction from dense multi-view input is largely solved.
Feed-forward generalisation to sparse-view and single-image input has
seen dramatic advances in the past two years.
And the first generation of world foundation models trained at internet scale
has demonstrated that video generation and spatial synthesis are convergent problems.
Yet the capability table in Section~\ref{sec:world} (Table~\ref{tab:nvs_wm})
reveals that three of the five world-modeling desiderata---temporal dynamics
(D3), action conditioning (D4), and physical plausibility (D5)---are only
partially satisfied by current best methods.
This section identifies the most important open problems that must be resolved
to cross these frontiers, and proposes a concrete research roadmap.

\subsection{The Geometry--Generation Trade-Off}

The fundamental tension between the \emph{reconstruction} and \emph{generative}
paradigms has not yet been resolved.
Reconstruction-based methods (NeRF, 3DGS, feed-forward Gaussian) deliver
geometric precision but are limited to what has been observed: they cannot
hallucinate content outside the observed view frustum.
Generative methods (diffusion NVS, world models) provide rich priors over
plausible scene completions but can violate multi-view geometric consistency,
producing views that are individually realistic but mutually contradictory in
3-D space.

Methods like ReconFusion~\citep{wu2024reconfusion},
Cat3D~\citep{gao2024cat3d}, and ViewCrafter~\citep{yu2024viewcrafter}
begin to bridge this gap by coupling a generative prior to an explicit 3-D
scaffold, but the integration is primarily architectural (conditioning the
diffusion model on a point cloud or NeRF) rather than principled (ensuring
that the generated content satisfies 3-D consistency constraints during
denoising).
Principled solutions---such as geometry-grounded diffusion that enforces
epipolar constraints during denoising, or diffusion-seeded 3DGS that
initialises the Gaussian field from a generated multi-view set---represent
an active frontier with significant room for innovation.
The key insight we advocate is that the 3-D consistency constraint should
be enforced \emph{during} generation rather than as a post-hoc filter,
which would require integrating differentiable rendering into the
diffusion denoising loop.

\subsection{Long-Horizon Temporal Coherence}

Existing 4-D methods maintain coherence over a few seconds, corresponding
to a few hundred frames of video.
World models for autonomous driving and robotics require coherence over minutes
of operation, during which the scene may change substantially due to agent
interactions, illumination changes, and new scene content entering the field of view.
Two promising directions have been identified in the recent literature,
though neither has yet achieved the long-horizon coherence required by
real-world applications.

The first direction is \emph{persistent scene memory}: Gaussian maps that
accumulate observations across time and preserve global consistency by storing
the entire world state as a growing set of Gaussian primitives.
Street Gaussians~\citep{yan2024streetgaussian} and OmniRe~\citep{chen2024omnire}
demonstrate that this approach can handle scene-length (several-minute) driving
sequences, but their memory grows without bound and they do not support
closed-loop rollouts of novel trajectories.
An efficient, bounded-memory Gaussian world state that can be queried and
updated in real time remains an open problem.

The second direction is \emph{world model pretraining}: leveraging internet-scale
video prediction as a temporal prior that can be fine-tuned to specific domains.
Cosmos~\citep{agarwal2025cosmos} and Genie~2~\citep{parkerholder2024genie2}
represent the current frontier here, but their coherence degrades over longer
horizons due to accumulated prediction errors in the autoregressive generation
process---a fundamental challenge that may require non-autoregressive or
hierarchical temporal architectures.

\subsection{Physical Grounding}

Current methods learn the appearance correlations that arise from physical
laws, but they do not represent physical laws explicitly.
This distinction matters for world modeling: a model that has learned that
objects tend to remain stationary may generate plausible-looking video, but
it will fail to predict the correct response to an unusual physical interaction
(pushing an object off a table, flooding a scene, deforming an elastic material).
Physical grounding is therefore a prerequisite for reliable world models in
applications that require reasoning about novel or unusual physical events.

PAC-NeRF~\citep{li2023pac} demonstrates a promising direction by integrating
continuum mechanics into NeRF, enabling system identification of material
properties from observed deformation.
PhysGaussian~\citep{xie2024physgaussian} brings physics simulation to Gaussian
Splatting, embedding Material Point Method (MPM) simulation directly in the
Gaussian optimisation loop.
The central challenge is extending these approaches to complex, open-world
scenes with diverse materials and contact interactions: both PAC-NeRF and
PhysGaussian have been demonstrated primarily on single-material, simple-geometry
scenarios.
Integrating differentiable physics more generally into the rendering pipeline
in a way that scales to complex scenes and diverse material properties remains
a largely open problem.

\subsection{Open-Vocabulary Semantic Integration}

The methods surveyed in this paper largely address appearance-level scene
understanding, but a full world model must also reason about \emph{semantics}:
what is this object, what can I do with it, how will it respond to my action?
The recent language-grounded 3-D field represents the most direct progress
toward this goal.
LERF~\citep{kerr2023lerf} grounds CLIP features in 3-D via NeRF,
enabling open-vocabulary 3-D queries such as ``where is the wooden chair?''.
LangSplat~\citep{qin2024langsplat} brings language grounding to 3DGS,
achieving real-time semantic query response while maintaining the rendering
speed advantage of the Gaussian representation.
Feature3DGS~\citep{zhou2024feature3dgs} provides a more general framework,
distilling arbitrary 2-D foundation model features---including SAM~\citep{kirillov2023sam}
and DINO---into the Gaussian representation.

Extending these language-grounding approaches to dynamic, actionable scenes
is a key step toward language-grounded world models: a model that can both
generate novel views of a scene and answer natural language queries about
what it contains, how it will evolve, and how an agent can interact with it.
This integration of spatial and semantic reasoning is a central challenge for
the next generation of world models.

\subsection{Evaluation and Benchmarks}

The evaluation gap identified in Section~\ref{sec:eval} requires a new
benchmark design that simultaneously tests all five world-modeling desiderata.
We propose a \textbf{World Modeling Evaluation Suite} with five axes:
\begin{enumerate}[leftmargin=*,label=(\roman*)]
  \item \emph{3-D consistency} measured by multi-view photometric consistency
    under adversarial viewpoints that maximally stress the geometric model;
  \item \emph{Temporal coherence} measured over long horizons ($>$30\,s) using
    feature-space consistency and structural similarity;
  \item \emph{Action responsiveness} measured in interactive settings by
    the quality of generated responses to specified agent actions;
  \item \emph{Physical plausibility} verified by comparison against
    physics-engine simulation of the same scenario;
  \item \emph{Open-world generalisation} measured by single-image synthesis
    quality across highly diverse, out-of-distribution environments.
\end{enumerate}
No existing benchmark covers all five axes.
We note that designing such a benchmark is itself a substantial research
contribution, as it requires defining a controlled test-bed for all five
dimensions without introducing confounding factors between them.

\subsection{Research Roadmap}

Drawing the threads of this analysis together, we propose five research
directions as the most productive frontiers for the immediate future.

\begin{enumerate}[leftmargin=*,label=\textbf{R\arabic*.}]
  \item \textbf{Unified geometry-generation architectures.}
    Develop architectures that couple explicit 3-D representations
    (NeRF, 3DGS) with diffusion priors in a manner that enforces
    3-D consistency during generation, not merely as a post-hoc filter.
    The key technical challenge is integrating differentiable rendering into
    the denoising loop without prohibitive computational cost.

  \item \textbf{Internet-scale feed-forward generalisation.}
    Scale feed-forward reconstruction models (LRM, DUSt3R, VGGT) to
    internet-scale pretraining on diverse unstructured video, enabling
    open-world NVS from a single image with generalisation to arbitrary
    scene types encountered at deployment time.

  \item \textbf{Differentiable physics in 4-D world models.}
    Integrate differentiable physics priors into 4-D world model training,
    building on the PAC-NeRF and PhysGaussian foundations, to enable
    physically grounded prediction of novel physical interactions.

  \item \textbf{Language and semantic grounding for action conditioning.}
    Extend language-grounded 3-D representations (LERF, LangSplat, Feature3DGS)
    to dynamic, actionable world models that can respond to natural language
    action specifications, enabling instruction-following spatial agents.

  \item \textbf{A comprehensive world modeling evaluation suite.}
    Design and release a benchmark covering all five world-modeling
    desiderata, with controlled, reproducible evaluation protocols
    that separate each axis's contribution to total system performance.
    This benchmark is a prerequisite for measuring progress on the
    most challenging open problems.
\end{enumerate}

Together, these five directions define a research agenda that, if successfully
prosecuted, would complete the transition from current \nvs{} systems---which
solve D1 and partially D2---to full world models that satisfy all five desiderata.

% ─────────────────────────────────────────────────────────────────────────────
